% Part: first-order-logic
% Chapter: beyond
% Section: intuitionistic-logic

\documentclass[../../../include/open-logic-section]{subfiles}

\begin{document}

\olfileid{fol}{byd}{il}

\olsection{স্বজ্ঞাবাদী যুক্তিবিদ্যা}

দ্বিতীয়-ক্রমের ও উচ্চতর-ক্রমের যুক্তিবিদ্যার বিপরীতে স্বজ্ঞাবাদী
প্রথম-ক্রমের যুক্তিবিদ্যা ধ্রুপদি সংস্করণটির একটি সীমিত রূপ; আরও
“নির্মাণমূলক” এক ধরনের যুক্তিবিচারের মডেল হওয়াই এর উদ্দেশ্য। নিচের
উদাহরণগুলি অন্তর্নিহিত কয়েকটি প্রেরণা বোঝাতে পারে।

ধরো, একদিন কেউ তোমার কাছে এসে ঘোষণা করল যে সে এমন একটি স্বাভাবিক
সংখ্যা~$x$ নির্ণয় করেছে যার ধর্ম হলো: $x$ মৌলিক হলে রিমান অনুমানটি
সত্য, আর $x$ যৌগিক হলে রিমান অনুমানটি মিথ্যা। কী দারুণ খবর!{} রিমান
অনুমান সত্য কি না, তা গণিতের বড় অমীমাংসিত প্রশ্নগুলির একটি; আর এখানে
মনে হচ্ছে সমস্যাটিকে একটি গণনার প্রশ্নে নামিয়ে আনা হয়েছে---অর্থাৎ,
একটি নির্দিষ্ট সংখ্যা মৌলিক কি না তা নির্ণয় করলেই হবে।

$x$-এর সেই জাদুকরি মানটি কী? লোকটি এভাবে বর্ণনা করল: রিমান অনুমান
সত্য হলে $x$ হলো $7$, আর তা না হলে $9$।

রেগে গিয়ে তুমি টাকা ফেরত চাইলে। ধ্রুপদি দৃষ্টিভঙ্গি থেকে উপরের
বর্ণনাটি সত্যিই $x$-এর একটি একক মান নির্ধারণ করে; কিন্তু তুমি আসলে
চাও $x$-এর মানটি \emph{স্পষ্টভাবে} দেওয়া হোক।

আরেকটি, সম্ভবত কম কৃত্রিম উদাহরণ নেওয়া যাক। আমরা জানি, একটি অমূলদ
সংখ্যাকে মূলদ ঘাতে তুলে মূলদ ফল পাওয়া সম্ভব। যেমন,
$\sqrt{2}^2 = 2$। যেটি কম পরিষ্কার তা হলো, কোনো অমূলদ সংখ্যাকে
\emph{অমূলদ} ঘাতে তুলেও মূলদ ফল পাওয়া সম্ভব কি না। নিচের উপপাদ্যটি
ইতিবাচক উত্তর দেয়:

\begin{thm}
এমন অমূলদ সংখ্যা $a$ ও $b$ আছে যেন $a^b$ মূলদ।
\end{thm}

\begin{proof}
$\sqrt{2}^{\sqrt{2}}$ বিবেচনা করো। এটি মূলদ হলে কাজ শেষ: $a = b =
\sqrt{2}$ নেওয়া যায়। তা না হলে এটি অমূলদ। তখন
\[
(\sqrt{2}^{\sqrt{2}})^{\sqrt{2}} = \sqrt{2}^{\sqrt{2} \cdot
  \sqrt{2}} = \sqrt{2}^2 = 2,
\]
যা অবশ্যই মূলদ। তাই এই ক্ষেত্রে $a$ হিসেবে
$\sqrt{2}^{\sqrt{2}}$ এবং $b$ হিসেবে~$\sqrt 2$ নাও।
\end{proof}

এটি কি বৈধ প্রমাণ? অধিকাংশ গণিতবিদ মনে করেন, হ্যাঁ। কিন্তু এখানেও
একটু অসন্তোষ থেকে যায়: একটি নির্দিষ্ট ধর্মবিশিষ্ট বাস্তব সংখ্যার
একটি জোড়ার অস্তিত্ব আমরা প্রমাণ করেছি, অথচ সেটি \emph{কোন} জোড়া তা
বলতে পারিনি। একই ফল এমনভাবে প্রমাণ করা সম্ভব যাতে প্রমাণের মধ্যেই
$a$, $b$ জোড়াটি \emph{দেওয়া} থাকে: $a = \sqrt{3}$ এবং $b = \log_3
4$ নাও। তাহলে
\[
a^b = \sqrt{3}^{\log_3 4} = 3^{1/2 \cdot \log_3 4} = (3^{\log_3
  4})^{1/2} = 4^{1/2}= 2,
\]
কারণ $3^{\log_3 x} = x$।

প্রথম প্রমাণটির মতো পদক্ষেপ নিষিদ্ধ এমন এক ধরনের যুক্তিবিচারের মডেল
হিসেবে স্বজ্ঞাবাদী যুক্তিবিদ্যা নির্মিত। $!A(x)$ পরিতৃপ্তকারী কোনো
$x$-এর অস্তিত্ব প্রমাণ করার অর্থ হলো, দ্বিতীয় প্রমাণের মতো একটি
নির্দিষ্ট~$x$ এবং সেটি $!A$ পরিতৃপ্ত করে তার প্রমাণ দিতে হবে। $!A$
অথবা $!B$ সত্য প্রমাণ করার জন্য দুটির কোনো একটিকে প্রমাণ করতে পারতে
হবে।

আনুষ্ঠানিকভাবে বলতে গেলে, প্রথম-ক্রমের যুক্তিবিদ্যার !!a{derivation}
ব্যবস্থাকে একটি নির্দিষ্ট উপায়ে সীমিত করলে স্বজ্ঞাবাদী প্রথম-ক্রমের
যুক্তিবিদ্যা পাওয়া যায়। একইভাবে দ্বিতীয়-ক্রমের বা উচ্চতর-ক্রমের
যুক্তিবিদ্যারও স্বজ্ঞাবাদী সংস্করণ আছে। গাণিতিক দৃষ্টিতে এগুলি নিছক
আনুষ্ঠানিক অবরোহী ব্যবস্থা; কিন্তু আগেই বলা হয়েছে, এক ধরনের গাণিতিক
যুক্তিবিচারের মডেল হওয়াই এগুলির উদ্দেশ্য। কেউ একে গণিত-সম্বন্ধীয়
কোনো নির্দিষ্ট দার্শনিক অবস্থানে (যেমন ব্রাউয়ারের স্বজ্ঞাবাদে) ন্যায্য
যুক্তিবিচার হিসেবে নিতে পারেন; কেউ একে আরও “মূর্ত” ও সন্তোষজনক এক
ধরনের গাণিতিক যুক্তিবিচার হিসেবে নিতে পারেন (বিশপের নির্মাণবাদের
ধারায়); আবার আনুষ্ঠানিক বর্ণনাটি অনানুষ্ঠানিক প্রেরণাকে ঠিকভাবে ধারণ
করে কি না, তা নিয়ে তর্কও করতে পারেন। কিন্তু আমাদের দার্শনিক অবস্থান
যা-ই হোক, আনুষ্ঠানিকভাবে উপস্থাপিত একটি যুক্তিবিদ্যা হিসেবে
স্বজ্ঞাবাদী যুক্তিবিদ্যাকে অধ্যয়ন করা যায়; এবং যে কারণেই হোক, অনেক
গাণিতিক যুক্তিবিদ তা করতে আগ্রহী।

স্বজ্ঞাবাদী সংযোজকগুলির একটি অনানুষ্ঠানিক নির্মাণমূলক ব্যাখ্যা আছে,
যা সাধারণত BHK ব্যাখ্যা নামে পরিচিত (ব্রাউয়ার, হেইটিং ও কলমগরভের
নামানুসারে)। ব্যাখ্যাটি এমন: $!A \land !B$-এর প্রমাণ গঠিত হয় $!A$-এর
একটি প্রমাণের সঙ্গে $!B$-এর একটি প্রমাণ জোড়া দিয়ে; $!A \lor !B$-এর
প্রমাণ গঠিত হয় হয় $!A$-এর একটি প্রমাণ, নয় $!B$-এর একটি প্রমাণ দিয়ে,
এবং কোনটি ঘটেছে তার স্পষ্ট তথ্যও আমাদের থাকে; $!A \lif !B$-এর প্রমাণ
গঠিত হয় এমন একটি প্রণালি দিয়ে যা $!A$-এর প্রমাণকে $!B$-এর প্রমাণে
রূপান্তর করে; $\lforall[x][!A(x)]$-এর প্রমাণ গঠিত হয় এমন একটি প্রণালি
দিয়ে যা~$x$-এর যেকোনো মানের জন্য $!A(x)$-এর একটি প্রমাণ ফেরায়; আর
$\lexists[x][!A(x)]$-এর প্রমাণ গঠিত হয়~$x$-এর একটি মান এবং সেই মানটি
$!A$ পরিতৃপ্ত করে তার প্রমাণ দিয়ে। “কারি--হাওয়ার্ড সমরূপতা” বা
“!!{formula}s-কে-টাইপ দৃষ্টান্ত” নামে পরিচিত গণনামূলক ভাষায়ও এই
ব্যাখ্যা বর্ণনা করা যায়: !!a{formula}-কে একটি নির্দিষ্ট ধরনের
ডেটা-টাইপের নির্দিষ্টকরণ এবং প্রমাণকে সেই ডেটা-টাইপের গণনামূলক বস্তু
হিসেবে ভাবো; এই বস্তুগুলিই সংশ্লিষ্ট !!{formula} সত্য বলে দেখতে দেয়।

স্বজ্ঞাবাদী যুক্তিবিদ্যাকে প্রায়ই ধ্রুপদি যুক্তিবিদ্যা থেকে
বর্জিত-মধ্যম নীতি “বাদ দেওয়া” ব্যবস্থা হিসেবে ভাবা হয়। নিচের
উপপাদ্যটি কথাটিকে আরও নির্ভুল করে।
\begin{thm}
স্বজ্ঞাবাদীভাবে নিচের স্বতঃসিদ্ধ-ছকগুলি পরস্পর সমতুল্য:
\begin{enumerate}
\item $(\lnot !A \lif \lfalse) \lif !A$।
\item $!A \lor \lnot !A$
\item $\lnot \lnot !A \lif !A$
\end{enumerate}
\end{thm}
যেকোনো একটি ছক থেকে অন্য দুটির নিদর্শন পাওয়া স্বজ্ঞাবাদী
যুক্তিবিদ্যার একটি ভালো অনুশীলন।

ব্রাউয়ারের স্বজ্ঞাবাদের আনুষ্ঠানিক রূপ হিসেবে প্রস্তাবিত
স্বজ্ঞাবাদী প্রস্তাবনামূলক যুক্তিবিদ্যার প্রথম অবরোহী ব্যবস্থাগুলি
স্বাধীনভাবে কলমগরভ, গ্লিভেঙ্কো ও হেইটিং দিয়েছিলেন। স্বজ্ঞাবাদী
প্রথম-ক্রমের যুক্তিবিদ্যার (এবং স্বজ্ঞাবাদী গণিতের কিছু অংশের) প্রথম
আনুষ্ঠানিক রূপ হেইটিংয়ের। ধ্রুপদিভাবে বৈধ কয়েকটি ছক
স্বজ্ঞাবাদীভাবে বৈধ না হলেও অনেকগুলিই বৈধ।

\emph{দ্বিনঞর্থকতা অনুবাদ} ধ্রুপদি ও স্বজ্ঞাবাদী যুক্তিবিদ্যার মধ্যে
একটি গুরুত্বপূর্ণ সম্পর্ক বর্ণনা করে। নিচের মতো করে এটি আরোহক্রমে
সংজ্ঞায়িত (ধ্রুপদি !!{formula}~$!A$-এর “স্বজ্ঞাবাদী” অনুবাদ হিসেবে
$!A^N$-কে ভাবো):
\begin{align*}
!A^N & \ident \lnot\lnot !A \quad \text{পরমাণু !!{formula}s $!A$-এর জন্য} \\
(!A \land !B)^N & \ident (!A^N \land !B^N) \\
(!A \lor !B)^N & \ident  \lnot\lnot (!A^N \lor !B^N) \\
(!A \lif !B)^N & \ident (!A^N \lif !B^N) \\
(\lforall[x][!A])^N & \ident \lforall[x][!A^N] \\
(\lexists[x][!A])^N & \ident \lnot\lnot\lexists[x][!A^N]
\end{align*}
প্রস্তাবনামূলক যুক্তিবিদ্যার জন্য কলমগরভ ও গ্লিভেঙ্কোর এই অনুবাদের
সংস্করণ ছিল; বিধেয় যুক্তিবিদ্যার ক্ষেত্রে গ্যোডেল ও গেন্ৎসেন
স্বাধীনভাবে এটি দিয়েছিলেন। আমাদের আছে:

\begin{thm}
\begin{enumerate}
\item $!A \liff !A^N$ ধ্রুপদিভাবে প্রমাণযোগ্য।
\item $!A$ ধ্রুপদিভাবে প্রমাণযোগ্য হলে $!A^N$ স্বজ্ঞাবাদীভাবে
  প্রমাণযোগ্য।
\end{enumerate}
\end{thm}

এখন নিচের সংলাপটি কল্পনা করা যায়। ধ্রুপদি গণিতবিদ: “আমি $!A$ প্রমাণ
করেছি!” স্বজ্ঞাবাদী গণিতবিদ: “তোমার প্রমাণ বৈধ নয়। আসলে তুমি
$!A^N$ প্রমাণ করেছ।” ধ্রুপদি গণিতবিদ: “তাতেই আমি সন্তুষ্ট!” ধ্রুপদি
গণিতবিদের দৃষ্টিতে স্বজ্ঞাবাদী গণিতবিদ নিছক খুঁতখুঁত করছেন, কারণ
দুটি সমতুল্য। কিন্তু স্বজ্ঞাবাদী জোর দিয়ে বলেন, পার্থক্য আছে।

লক্ষ করো, উপরের অনুবাদটি কেবল বিশুদ্ধ যুক্তিবিদ্যা নিয়ে; ধ্রুপদি ও
স্বজ্ঞাবাদী গণিতের উপযুক্ত \emph{অ-যৌক্তিক} স্বতঃসিদ্ধগুলি কী, অথবা
তাদের পারস্পরিক সম্পর্ক কী, এই প্রশ্ন তার বিষয় নয়। তবে উপরের
উপপাদ্যটির নিচের সামান্য সম্প্রসারণ কিছু উপকারী তথ্য দেয়:

\begin{thm}
$\Gamma$ ধ্রুপদিভাবে $!A$ প্রমাণ করলে $\Gamma^N$ স্বজ্ঞাবাদীভাবে
$!A^N$ প্রমাণ করে।
\end{thm}

অন্যভাবে বললে, কয়েকটি অনুমান থেকে $!A$ ধ্রুপদিভাবে প্রমাণযোগ্য হলে
সেগুলির দ্বিনঞর্থকতা অনুবাদ থেকে $!A^N$ প্রমাণযোগ্য।

কোনো বাক্য বা প্রস্তাবনামূলক !!{formula} স্বজ্ঞাবাদীভাবে বৈধ দেখাতে
শুধু একটি প্রমাণ দিতে হয়। কিন্তু সেটি বৈধ নয় কীভাবে দেখাবে? এই
উদ্দেশ্যে এমন একটি অর্থতত্ত্ব দরকার যা নির্ভুল এবং সম্ভব হলে সম্পূর্ণ।
ক্রিপকে-প্রদত্ত একটি অর্থতত্ত্ব সুন্দরভাবে এই কাজ করে।

ধ্রুপদি যুক্তিবিদ্যার জন্য যে খেলা খেলেছিলাম, এখানেও তাই করা যায়:
অর্থতত্ত্ব সংজ্ঞায়িত করে নির্ভুলতা ও সম্পূর্ণতা প্রমাণ করা। তবে নিচের
পার্থক্যটি লক্ষ করার মতো। ধ্রুপদি যুক্তিবিদ্যার ক্ষেত্রে অর্থতত্ত্বটি
এক অর্থে সংযোজকগুলির অর্থের মধ্যেই নিহিত “স্বাভাবিক” অর্থতত্ত্ব ছিল।
ক্রিপকে অর্থতত্ত্বের কিছু স্বজ্ঞাত প্রেরণা দেওয়া গেলেও সেটিকে একইভাবে
অনিবার্য মনে হয় না। তাছাড়া, ধ্রুপদি !!{structure}-এর ধারণা একটি
স্বাভাবিক গাণিতিক ধারণা; তাই !!a{structure}-এর ধারণাকে ধ্রুপদি
প্রথম-ক্রমের যুক্তিবিদ্যা অধ্যয়নের উপকরণ হিসেবে নেওয়া যায়, আবার
ধ্রুপদি প্রথম-ক্রমের যুক্তিবিদ্যাকে গাণিতিক !!{structure}s অধ্যয়নের
উপকরণ হিসেবেও নেওয়া যায়। বিপরীতে, ক্রিপকে !!{structure}s-কে কেবল
যৌক্তিক নির্মাণ হিসেবেই দেখা যায়; এগুলির স্বাধীন গাণিতিক আগ্রহ আছে
বলে মনে হয় না।

কোনো প্রস্তাবনামূলক ভাষার জন্য একটি ক্রিপকে !!{structure}~$\mModel M
= \tuple{W, R, V}$ গঠিত হয় একটি সেট~$W$, $W$-এর উপর ক্ষুদ্রতম
!!{element}-সহ একটি আংশিক ক্রম~$R$, এবং $W$-এর !!{element}s-এ
প্রস্তাবনামূলক চলরাশিগুলির একটি “একঘেয়ে” আরোপ দিয়ে। স্বজ্ঞাটি হলো,
$W$-এর !!{element}s “জগৎ” বা “জ্ঞানের অবস্থা” নির্দেশ করে; $v \geq
u$ উপাদানটি~$u$-র একটি “সম্ভাব্য ভবিষ্যৎ অবস্থা” নির্দেশ করে; আর
প্রস্তাবনামূলক চলরাশি হিসেবে~$u$-তে আরোপিত প্রস্তাবনাগুলিই $u$ অবস্থায়
সত্য বলে জানা। বাধ্যকরণ সম্পর্ক $\mSat{M}{!A}[w]$ এই সম্পর্কটিকে
ভাষার যেকোনো !!{formula}s-এ প্রসারিত করে; $\mSat{M}{!A}[w]$-কে
“$!A$, $w$ অবস্থায় সত্য” বলে পড়ো। সম্পর্কটি নিচের মতো আরোহক্রমে
সংজ্ঞায়িত:
\begin{enumerate}
\item $\mSat{M}{\Obj p_i}[w]$ হবে যদি এবং কেবল যদি $\Obj p_i$ হলো
  $w$-তে আরোপিত প্রস্তাবনামূলক চলরাশিগুলির একটি।
\item $\mSat/{M}{\lfalse}[w]$।
\item $\mSat{M}{(!A \land !B)}[w]$ হবে যদি এবং কেবল যদি
  $\mSat{M}{!A}[w]$ এবং $\mSat{M}{!B}[w]$।
\item $\mSat{M}{(!A \lor !B)}[w]$ হবে যদি এবং কেবল যদি
  $\mSat{M}{!A}[w]$ অথবা $\mSat{M}{!B}[w]$।
\item $\mSat{M}{(!A \lif !B)}[w]$ হবে যদি এবং কেবল যদি, যখনই
  $w' \geq w$ এবং $\mSat{M}{!A}[w']$, তখন $\mSat{M}{!B}[w']$।
\end{enumerate}
$\lnot (p \land q) \lif (\lnot p \lor \lnot q)$ স্বজ্ঞাবাদীভাবে
বৈধ নয় দেখানোর চেষ্টা একটি ভালো অনুশীলন: একটি ক্রিপকে
!!{structure} বানাও যা তার একটি প্রতিমডেল দেয়।

\end{document}
